% -----------------------------------------------------------------------------
% Physical photon polarization projector after Maxwell quantization
% -----------------------------------------------------------------------------

Maxwell 约束把实光子的物理自由度限制为两个横向偏振。散射振幅中还需要把这两个偏振的求和写成协变张量。设实光子的物理四动量为 $k^\mu$，满足 $k^2=0$；再选一个满足 $k\!\cdot n\neq0$ 的辅助四矢量 $n^\mu$ 来固定偏振代表元。取
\begin{equation}
 k\!\cdot\epsilon_\lambda=0,
 \qquad
 n\!\cdot\epsilon_\lambda=0,
 \qquad
 \epsilon_\lambda\!\cdot\epsilon_{\lambda'}^*=-\delta_{\lambda\lambda'},
 \qquad \lambda=1,2.
 \label{eq:physical-photon-polarization-constraints-r68}
\end{equation}
两个物理偏振的完备张量为
\begin{equation}
 \boxed{
 \sum_{\lambda=1}^{2}
 \epsilon_\lambda^\mu(k)\epsilon_\lambda^{\nu *}(k)
 =-\eta^{\mu\nu}
 +\frac{k^\mu n^\nu+n^\mu k^\nu}{k\!\cdot n}
 -\frac{n^2k^\mu k^\nu}{(k\!\cdot n)^2}.}
 \label{eq:physical-photon-polarization-projector-dp42}
\end{equation}
辅助矢量 $n^\mu$ 只选择规范代表元，不进入规范不变的可观测量。若光子指标与守恒流 $J_\mu$ 收缩，$k^\mu J_\mu=0$ 使所有含 $k^\mu$ 的项消失，因此
\begin{equation}
 \boxed{
 J_\mu
 \left(\sum_{\lambda=1}^{2}
 \epsilon_\lambda^\mu\epsilon_\lambda^{\nu *}\right)
 J_\nu^*
 =-J_\mu J^{\mu *}.}
 \label{eq:physical-photon-polarization-conserved-current-dp42}
\end{equation}
所以常用的 $-\eta^{\mu\nu}$ 替换只在与守恒振幅收缩以后成立；它不是两个物理偏振子空间在完整四维 Minkowski 空间中的恒等算符。

\begin{exbox}[$z$ 方向光子的偏振与螺旋度]
取
\begin{equation*}
 k^\mu=\left(\frac{\hbar\omega}{c},0,0,\frac{\hbar\omega}{c}\right),
 \qquad k^2=0,
\end{equation*}
并在辐射规范代表元中取两个线偏振
\begin{equation*}
 \epsilon_x^\mu=(0,1,0,0),
 \qquad
 \epsilon_y^\mu=(0,0,1,0).
\end{equation*}
逐项可检验
\begin{equation*}
 k\!\cdot\epsilon_x=k\!\cdot\epsilon_y=0,
 \qquad
 \epsilon_x^2=\epsilon_y^2=-1,
 \qquad
 \epsilon_x\!\cdot\epsilon_y=0.
\end{equation*}
因此这两个四矢量确实张成与传播方向垂直的二维物理偏振空间。它们的显式完备和为
\begin{equation*}
 \epsilon_x^\mu\epsilon_x^\nu
 +\epsilon_y^\mu\epsilon_y^\nu
 =\operatorname{diag}(0,1,1,0),
\end{equation*}
这正是\eqnrefs{eq:physical-photon-polarization-projector-dp42} 在 $n^\mu=(1,0,0,0)$ 下的结果。

对角化绕传播轴的转动，更自然的基是圆偏振
\begin{equation}
 \epsilon_\lambda^\mu
 =\frac1{\sqrt2}
 \left(0,1,\ii\lambda,0\right),
 \qquad \lambda=\pm1.
 \label{eq:photon-circular-polarizations-dp62}
\end{equation}
三维矢量表示的 $z$ 轴转动生成元满足
\begin{equation*}
 (S_z)_{ij}=-\ii\epsilon_{3ij},
 \qquad
 S_z\,\boldsymbol\epsilon_\lambda
 =\lambda\boldsymbol\epsilon_\lambda.
\end{equation*}
所以有限转动给出
\begin{equation*}
 \ee^{-\ii\phi S_z}\boldsymbol\epsilon_\lambda
 =\ee^{-\ii\lambda\phi}\boldsymbol\epsilon_\lambda.
\end{equation*}
这说明 $\lambda=\pm1$ 正是沿 $\mathbf k$ 的两个光子螺旋度。线偏振与圆偏振并不是两类不同光子，而只是同一个二维物理偏振空间的两组基：
\begin{equation*}
 \epsilon_x=\frac{\epsilon_++\epsilon_-}{\sqrt2},
 \qquad
 \epsilon_y=\frac{\epsilon_+-\epsilon_-}{\ii\sqrt2}.
\end{equation*}
散射实验若不分辨光子偏振，求和的是这个二维空间的一组完备基，因此用线偏振基或螺旋度基计算必须给出相同的未偏振截面。
\end{exbox}


\begin{exbox}[有质量矢量场为什么多一个纵向偏振]
为了和光子比较，考虑质量为 $M>0$ 的自由 Proca 模，取四动量
\begin{equation*}
 k^\mu=\left(\frac Ec,0,0,p\right),
 \qquad
 k^2=M^2c^2,
 \qquad
 E^2=M^2c^4+c^2p^2.
\end{equation*}
两个横向偏振仍可取 $\epsilon_x^\mu$、$\epsilon_y^\mu$。但此时还存在第三个满足 $k\cdot\epsilon_L=0$ 的归一化纵向偏振：
\begin{equation}
 \boxed{
 \epsilon_L^\mu(k)
 =\left(\frac{p}{Mc},0,0,\frac{E}{Mc^2}\right).}
 \label{eq:massive-vector-longitudinal-polarization-dp62}
\end{equation}
逐项检查得到
\begin{align*}
 k\cdot\epsilon_L
 &=\frac Ec\frac{p}{Mc}-p\frac{E}{Mc^2}=0,\\
 \epsilon_L^2
 &=\frac{p^2}{M^2c^2}-\frac{E^2}{M^2c^4}=-1.
\end{align*}
因此有质量自旋 $1$ 粒子具有三个物理偏振。三者的完备和为
\begin{equation}
 \boxed{
 \sum_{\lambda=1,2,L}
 \epsilon_\lambda^\mu(k)\epsilon_\lambda^{\nu *}(k)
 =-\eta^{\mu\nu}+\frac{k^\mu k^\nu}{M^2c^2}.}
 \label{eq:massive-vector-polarization-sum-dp62}
\end{equation}
由上式的显式分量 $\epsilon_L^0=p/(Mc)$ 与 $\boldsymbol\epsilon_L=(E/Mc^2)\hat{\mathbf p}$ 可见，在固定非零三动量 $p$ 下取 $M\to0$ 时，各分量均为 $O(M^{-1})$，因而有质量矢量场的纵向偏振不存在逐分量平滑的无质量极限。无质量规范场出现规范冗余；规范约束把纵向和时间样自由度从物理 Hilbert 空间中移除，只留下前一个算例中的两个螺旋度。标准模型的 $W^\pm,Z$ 因 Higgs 机制有质量，纵向偏振成为真实自由度；它与 Higgs 场中沿真空对称轨道的切向标量涨落之间的关系由破缺规范理论的质量矩阵和约束结构确定。
\end{exbox}
